\documentclass[10pt]{article}
\usepackage[T1]{fontenc}
\usepackage[utf8]{inputenc}
\usepackage{amsmath}
\usepackage{amssymb}
\usepackage[margin=1in]{geometry}
\usepackage{enumitem}
\usepackage{graphicx}
\usepackage{booktabs}
\usepackage{caption}
\usepackage{hyperref}

\title{COSC 421 Final Report}
\date{\today}

\begin{document}
\maketitle

\section*{1. Introduction to Metrics}

In order to determine the correlation between structural centrality of the major Canadian airports and their role, if any, in delay propagation, several centrality metrics and delay metrics were compared. For analyzing structural centrality:
\begin{enumerate}[label=\arabic*.]
    \item Degree (in, out, and total)
    \item Strength (in, out, total, and log)
    \item Betweenness Centrality
    \item Closeness Centrality
    \item PageRank
    \item Undirected Eigenvector Centrality
\end{enumerate}

For analyzing delay propagation outcomes:
\begin{enumerate}[label=\arabic*.]
    \item n\_pairs (number of origin-destination pairs that involve the airport in question)
    \item n\_severe\_prev (number of severe ``predecessor'' events that arrived at this airport)
    \item n\_propagated (number of events which propagate from this airport)
    \item propagation\_rate (proportion of severe predecessors that propagate)
    \item amplification (index of whether delays are generally amplified or dampened by this airport)
    \item rho\_spearman (Spearman correlation between previous and subsequent delays with this airport)
\end{enumerate}

\section*{2. Centrality and Delays}

\subsection*{2.1 Centrality Hierarchy}

By examining total strength, closeness centrality, PageRank, and eigenvector centrality, an idea of the structural layout of Canada's major airports becomes clear. YYZ, YYC, and YVR form the core, with YUL as an intermediary, and Edmonton at the comparative periphery of the network.

\begin{table}[h!]
    \centering
    \caption{Centrality metrics for major Canadian airports}
    \begin{tabular}{lcccc}
        \toprule
        Airport & strength\_total & closeness & eigenvector & PageRank \\
        \midrule
        YYZ (Toronto)   & 2342           & 255.64              & 1.00             & 0.260             \\
        YYC (Calgary)   & $\approx 2000$ & $\approx 235$--$251$ & $\approx 0.95$   & $\approx 0.237$--$0.238$ \\
        YVR (Vancouver) & $\approx 2000$ & $\approx 235$--$251$ & $\approx 0.95$   & $\approx 0.237$--$0.238$ \\
        YUL (Montreal)  & 1127           & 171                 & $\approx 0.594$  & $\approx 0.136$   \\
        YEG (Edmonton)  & 789            & 85.5                & $\approx 0.407$  & $\approx 0.128$   \\
        \bottomrule
    \end{tabular}
\end{table}

YYZ, YYC, and YVR all share high values in all measures of centrality. Comparatively, Montreal is significantly lower, with most metrics being about half of the top three. Edmonton is even less connected. Additionally, this data differentiates Toronto as a bridge. YYZ has a betweenness of 0.5, higher than that of YYC at 0.167, while the others have zero. 

\subsection*{2.2 Delay Volumes and Propagation}

\begin{table}[h!]
    \centering
    \caption{Delay volume and propagation metrics by airport}
    \begin{tabular}{lcccc}
        \toprule
        Airport & n\_severe\_prev & n\_propagated & propagation\_rate & n\_propagated / strength\_total \\
        \midrule
        YYZ & 124 & 26 & 0.210 & 0.0111 \\
        YYC & 114 & 28 & 0.246 & 0.0137 \\
        YVR & 94  & 28 & 0.298 & 0.0138 \\
        YUL & 54  & 7  & 0.130 & 0.0062 \\
        YEG & 23  & 6  & 0.261 & 0.0076 \\
        \bottomrule
    \end{tabular}
\end{table}

YYZ, YYC, and YVR, unquestionably the most central airports, also account for the majority of both severe incoming delay events and delay propagation. On a rate basis, Vancouver and Calgary are more likely to propagate severe events, both per event and per unit of traffic, when compared to Montreal and Edmonton, with Toronto slightly lower in propagation rate despite its high centrality.

\section*{3. Centrality and Delay Relationships}

Since the analysis focused on only Canada's five major airports, correlation magnitudes are useful more as a descriptive measure. However, they still reveal important patterns.

\subsection*{3.1 Centrality vs. Delay Volume}

For volume metrics (n\_pairs, n\_severe\_prev, n\_propagated), correlations with centrality are very high, as demonstrated by the Pearson correlation coefficient.

\[
r_{XY}
=
\frac{\displaystyle\sum_{i=1}^{n} (x_i - \bar{x})(y_i - \bar{y})}
     {\sqrt{\displaystyle\sum_{i=1}^{n} (x_i - \bar{x})^2}
      \sqrt{\displaystyle\sum_{i=1}^{n} (y_i - \bar{y})^2}}
\]

\begin{table}[h!]
    \centering
    \caption{Pearson correlations between centrality and delay volume metrics}
    \begin{tabular}{lcccc}
        \toprule
        Metric 1 & Metric 2 & Pearson \\
        \midrule
        strength\_total & n\_severe\_prev & $\approx 0.98$ \\
        strength\_total & n\_pairs        & $\approx 0.97$ \\
        strength\_total & n\_propagated   & $\approx 0.95$ \\
        closeness       & n\_pairs        & $\approx 0.998$ \\
        closeness       & n\_severe\_prev & $\approx 0.98$ \\
        pagerank        & n\_severe\_prev & $\approx 0.96$ \\
        eigen\_undirected & n\_propagated & $\approx 0.97$ \\
        \bottomrule
    \end{tabular}
\end{table}
\end{document}
